\section{Remarks}

\subsection{State Lists}

The current electoral procedure provides that the Bundestag mandates allocated to a party
are further divided among state lists. State lists would no longer exist under the procedure proposed here.
States should play no role in federal elections. Regional
representation already takes place through the proposed distribution of constituencies.
Each party running for election must run in the entire electoral area, i.e., in the federal election in the
entire federal territory.

The parties may internally maintain state lists if they wish, in order to internally carry out a
state balance in filling the Bundestag list. But that is their
decision. They have free disposal over the half of their parliamentary mandates not elected by citizens.

\subsection{Five Percent Threshold}

The presented procedure already contains a hurdle through the requirement to run in all constituencies with three candidates,
which could be a problem for small parties. The three
required lead candidates could also be a hurdle for smaller parties. In addition,
the \enquote{natural threshold} (only whole-number representatives) has an even stronger effect due to the
\enquote{compression} of the opposition.

For these reasons, I think one could possibly dispense with an explicit five percent threshold.
In the simulation, it is not considered. However, if it should nevertheless be politically desired,
it could easily be incorporated into the procedure.

\subsection{No More CDU/CSU Alliance}

From a state theory perspective, it is absolutely necessary that all parties and lead candidates running for election
run in the entire electoral area. Every parliamentary and
government member is obligated to the entire people and may not only represent the interests of a
specific group or a specific partial electoral area. Currently, the electoral success
of the CSU depends exclusively on how it is received by Bavarian voters and not on what the
rest of Germany thinks of it. As a result, we have state representation in the Bundestag through the back door
and the usual difficulties between coalition partners,
even if the CDU/CSU were to win the absolute majority. Governing means making impartial
trade-offs between different interests. Interest representation is lobbying.
Interest representation is very necessary and nothing bad in itself. The governing
can only consider interests in their deliberations if they also know them. However, the
interests must be represented before the impartial and not by the impartial.
Interest representation of the states already takes place in the Bundesrat. Bavaria is also
represented there. That must be enough. Why should Bavaria be more privileged than all other
federal states?

It is actually incomprehensible and can only be explained historically that the Federal Constitutional Court
has allowed a party running for the Bundestag not to run in all states.
As wrong as the CSU special arrangement already is now, it would be even more wrong under the
electoral procedure proposed here. The electoral procedure specifically ensures that coalitions remain
possible but are no longer necessary to be able to govern. Through this CDU/CSU combination, however, we have the coalition guaranteed, even with only an awarded absolute
majority. This would make the entire procedure absurd. If both CDU and CSU run in
all federal states, they can gladly form a coalition afterward if they want to. However, if the
CSU does not want to disappear into insignificance at the federal level, it must develop
concepts that are well received by all of Germany and not just by Bavaria.
Perhaps it will succeed in this, as apparently some things do run better in Bavaria than
in the rest of the republic. The state fiscal equalization speaks a clear language. The way to
roll it out to all of Germany is to take responsibility for all of Germany yourself by
running for election in all of Germany. Perhaps then it will also work with the Bavarian
Federal Chancellor!

\clearpage
\subsection{Multiple Christian Parties}

CDU/CSU could also merge into one party to meet the requirement for
parties running nationwide. This would be completely fine from a state theory perspective, and perhaps
they will indeed do so.

However, it would be desirable, though not absolutely necessary, for several large
Christian parties to run nationwide for election, so that convinced Christian voters also get a
real political choice and do not have to vote only according to religious affiliation. Currently,
many convinced Christians vote for CDU/CSU because of the \enquote{C} in the name, although they do not
necessarily agree with their political views. If several Christian parties with
different political views were available to choose from, this would certainly be beneficial for
democracy.
